% §7  Case Study: Graded Framing for Cache-Timing Side Channels
% Self-contained; \input into main POPL paper at §7.
% Assumes: \usepackage{amsmath,amssymb,algorithm,algpseudocode,booktabs,stmaryrd}

\section{Case Study: Cache-Timing Side Channels in CT-IMP}
\label{sec:case-study}

We instantiate the cocycle-graded frame rule on a concrete
cache-timing side-channel example, deriving computable, tight
defect bounds and recovering a quantitative noninterference
result as a corollary.

%------------------------------------------------------------------
\subsection{The CT-IMP Language}
\label{sec:ctimp}

\textbf{Syntax.}  CT-IMP is a small imperative language extended with
explicit cache-visible memory accesses:
\begin{align*}
  e &::= n \mid x \mid e_1 \mathop{\mathit{op}} e_2 \\
  c &::= \mathbf{skip} \mid x \mathbin{:=} e \mid c_1 \mathbin{;} c_2
       \mid \mathbf{if}\; e\; \mathbf{then}\; c_1\; \mathbf{else}\; c_2 \\
    &\quad\mid \mathbf{while}\; e\; \mathbf{do}\; c
       \mid \mathbf{cache\_load}(x, e)
       \mid \mathbf{store}(e, e')
\end{align*}
The instruction $\mathbf{cache\_load}(x, e)$ reads the word at address
$\mathit{eval}(e)$ into register $x$ \emph{and} updates the
observable cache state; $\mathbf{store}$ is a silent write.

\textbf{Machine state.}
A \emph{configuration} is a triple $\langle c,\sigma,\kappa\rangle$
where $\sigma : \mathit{Var} \to \mathit{Val}$ is a register store and
$\kappa : \mathit{CacheLine} \to \{\mathsf{hot},\mathsf{cold}\}$
records which $B = 64$-byte cache lines are currently resident.
The heap $h : \mathit{Loc} \rightharpoonup \mathit{Val}$ is implicit.

%------------------------------------------------------------------
\subsection{Small-Step Semantics}
\label{sec:semantics}

\begin{figure*}[t]
\small
\[
\frac{a = \mathit{eval}(\sigma, e) \quad \ell = \lfloor a/B\rfloor
      \quad \kappa(\ell) = \mathsf{hot} \quad \sigma' = \sigma[x \mapsto h(a)]}
     {\langle \mathbf{cache\_load}(x,e),\sigma,\kappa\rangle
      \xrightarrow{t_\mathsf{hit}}
      \langle \mathbf{skip},\sigma',\kappa\rangle}
\;\textsc{(Load-Hit)}
\]
\[
\frac{a = \mathit{eval}(\sigma, e) \quad \ell = \lfloor a/B\rfloor
      \quad \kappa(\ell) = \mathsf{cold}
      \quad \kappa' = \mathit{lru}(\kappa,\ell)
      \quad \sigma' = \sigma[x \mapsto h(a)]}
     {\langle \mathbf{cache\_load}(x,e),\sigma,\kappa\rangle
      \xrightarrow{t_\mathsf{miss}}
      \langle \mathbf{skip},\sigma',\kappa'\rangle}
\;\textsc{(Load-Miss)}
\]
\[
\frac{\langle c_1,\sigma,\kappa\rangle \xrightarrow{t}
      \langle c_1',\sigma',\kappa'\rangle}
     {\langle c_1;c_2,\sigma,\kappa\rangle \xrightarrow{t}
      \langle c_1';c_2,\sigma',\kappa'\rangle}
\;\textsc{(Seq)}
\qquad
\frac{}{\langle x\mathbin{:=}e,\sigma,\kappa\rangle \xrightarrow{1}
        \langle\mathbf{skip},\sigma[x\mapsto\mathit{eval}(\sigma,e)],\kappa\rangle}
\;\textsc{(Assign)}
\]
\caption{Selected small-step rules for CT-IMP.
  $t_\mathsf{hit}\!=\!1$\,ns, $t_\mathsf{miss}\!=\!100$\,ns,
  $B\!=\!64$\,bytes.
  $\mathit{lru}(\kappa,\ell)$ evicts the least-recently-used cold line
  and marks $\ell$ hot.
  The total execution time of a program is
  $\mathit{time}(c,\sigma,\kappa) = \sum_i t_i$ over all reduction steps.}
\label{fig:ctimp-rules}
\end{figure*}

The \emph{timing cost function} is
$\mathit{time}(c,\sigma,\kappa) \in \mathbb{N}$,
defined as the sum of step labels in the reduction sequence to
$\mathbf{skip}$.
We write $\mathit{fp}(c,h)$ for the \emph{cache footprint}:
the set of cache lines accessed during execution of $c$ from heap $h$.

%------------------------------------------------------------------
\subsection{Two Password-Checker Variants}

\begin{figure}[h]
\begin{algorithmic}[1]
\Procedure{check\_ct}{$\mathit{pwd}$, $\mathit{secret}$, $N$}
  \State $\mathit{match} \gets 0$;\; $i \gets 0$
  \While{$i < N$}
    \State $\mathbf{cache\_load}(s,\; \mathit{secret}+i)$
    \State $\mathbf{cache\_load}(t,\; \mathit{pwd}+i)$
    \State $\mathit{match} \gets \mathit{match} + (s = t)$
    \State $i \gets i + 1$
  \EndWhile
  \State \Return $(\mathit{match} = N)$
\EndProcedure
\end{algorithmic}
\smallskip
\begin{algorithmic}[1]
\Procedure{check\_vuln}{$\mathit{pwd}$, $\mathit{secret}$, $N$}
  \State $i \gets 0$
  \While{$i < N$}
    \State $\mathbf{cache\_load}(s,\; \mathit{secret}+i)$
    \State $\mathbf{cache\_load}(t,\; \mathit{pwd}+i)$
    \If{$s \neq t$} \Return $\mathbf{false}$ \EndIf  \Comment{early exit!}
    \State $i \gets i + 1$
  \EndWhile
  \State \Return $\mathbf{true}$
\EndProcedure
\end{algorithmic}
\caption{Constant-time (top) and vulnerable (bottom) password checkers.
Both compare an $N$-byte \emph{secret} to a user-supplied
\emph{pwd}.  The constant-time variant always performs $N$ loads from
each array; the vulnerable variant exits at the first mismatch.}
\label{fig:programs}
\end{figure}

%------------------------------------------------------------------
\subsection{Graded Triples}

\paragraph{Effect monoid and interaction measure.}
Take the effect monoid $M = (\mathcal{P}(\mathit{CacheLine}), \cup, \emptyset)$
where the grade of a command $c$ is $\mathit{fp}(c,h)$.
Define the interaction measure:
\begin{equation}
  r(S, D) = \bigl|\, S \cap D \,\bigr| \cdot (t_\mathsf{miss} - t_\mathsf{hit})
  \label{eq:interaction}
\end{equation}
so $r(S,D)$ counts how many cache lines in $D$ are shared with
(and thus potentially pre-warmed by) the computation $S$.
Right-linearity (RL) holds since $r(S, D_1 \cup D_2) = r(S,D_1) + r(S,D_2)$
when $D_1 \cap D_2 = \emptyset$.

The \emph{interference cocycle} from Definition~2.3 of the main paper is:
\begin{equation}
  c(S_1, S_2, D)
    = r(S_1 \cup S_2, D) - r(S_1, D) - r(S_2, D)
    = -\,|S_1 \cap S_2 \cap D| \cdot (t_\mathsf{miss} - t_\mathsf{hit})
  \label{eq:cocycle}
\end{equation}
Note $c \leq 0$ always: composing two computations sharing cache lines
with $D$ saves, never wastes, time relative to the additive estimate.
The norm $\|c(S_1,S_2,D)\|_\infty
= |S_1\cap S_2\cap D|\cdot(t_\mathsf{miss}-t_\mathsf{hit})$
is the \emph{cache-sharing bonus}, i.e., the frame-rule defect.

\paragraph{Constant-time variant.}
Let $h_s = \mathit{owns}(\mathit{secret}, N)$
and $h_p = \mathit{owns}(\mathit{pwd}, N)$.
Assume the allocator places \texttt{secret} and \texttt{pwd} in
\emph{disjoint} cache sets (a standard hardening assumption).
Then $\mathit{fp}(\mathtt{check\_ct},h_s) \cap
      \mathit{fp}(\mathtt{check\_ct},h_p) = \emptyset$, so
$c(S_s, S_p, D) = 0$ for all $D$.
The graded triple is:
\begin{equation}
  \vdash
  \bigl\{\,h_s * h_p\,\bigr\}\;\;
  \mathtt{check\_ct}
  \;\;\bigl\{\,\exists v.\, h_s * h_p * (\mathit{result}=v)\,\bigr\}
  \;_{\varepsilon=0}
  \label{eq:ct-triple}
\end{equation}
\emph{Clean frame}: the frame rule holds exactly; timing depends only
on $N$, not on the values of \texttt{secret}.

\paragraph{Vulnerable variant.}
Let $k \in \{0,\ldots,N-1\}$ be the index of the first mismatch
(geometrically distributed with parameter $p = 1/256$ for random
passwords over a byte alphabet).
On each iteration $i \leq k$, both loads execute; on iteration $k$
the early exit fires.
The expected number of cache misses in the \texttt{secret} footprint
is $k+1$; the early-exit leaks $k$ through timing.
Formally:
\begin{align}
  \mathit{time}(\mathtt{check\_vuln}, h_s, h_p)
    &= (k+1)\cdot(t_\mathsf{hit} + t_\mathsf{miss}) + O(1)
    \label{eq:vuln-time}
\end{align}
where the $O(1)$ covers branch costs.
The cocycle defect accumulated over the early-exit loop is:
\begin{align}
  \varepsilon_\mathit{vuln}
  &= \sum_{i=0}^{N-1}
     \Pr[\text{exit at step }i]\cdot i \cdot t_\mathsf{miss}
    + N\cdot\|c\|_\infty
  \notag\\
  &= \mathbb{E}[k]\cdot t_\mathsf{miss} + N\cdot\|c\|_\infty
  \notag\\
  &= \frac{1-p}{p}\cdot t_\mathsf{miss}
     + N\cdot(t_\mathsf{miss}-t_\mathsf{hit})
  \label{eq:defect}
\end{align}
For $p = 1/256$, $N = 16$, $t_\mathsf{hit}=1$\,ns,
$t_\mathsf{miss}=100$\,ns:
$\varepsilon_\mathit{vuln} = 255\cdot 100 + 16\cdot 99 = 27{,}084$\,ns.

The graded triple for the vulnerable variant is:
\begin{equation}
  \vdash
  \bigl\{\,h_s * h_p\,\bigr\}
  \;\;\mathtt{check\_vuln}\;\;
  \bigl\{\,\exists v.\, h_s * h_p * (\mathit{result}=v)\,\bigr\}
  \;_{\varepsilon_\mathit{vuln}}
  \label{eq:vuln-triple}
\end{equation}

%------------------------------------------------------------------
\subsection{Quantitative Noninterference as a Corollary}

\begin{theorem}[Timing Leakage Bound]
\label{thm:nonint}
Let $\mathit{secret}_0, \mathit{secret}_1 \in \{0,\ldots,255\}^N$
be any two secrets of the same length.
Under the graded triple~\eqref{eq:vuln-triple}:
\begin{equation}
  \sup_t \bigl|\Pr[\mathit{time}(\mathtt{check\_vuln},\mathit{secret}_0)=t]
         - \Pr[\mathit{time}(\mathtt{check\_vuln},\mathit{secret}_1)=t]\bigr|
  \;\leq\;
  \frac{\varepsilon_\mathit{vuln}}{t_\mathsf{miss} - t_\mathsf{hit}}
  \label{eq:nonint-bound}
\end{equation}
\end{theorem}
\begin{proof}[Proof sketch]
The cocycle-graded frame rule (Theorem~5.1) gives that the postcondition
timing deviates by at most $\varepsilon_\mathit{vuln}$ across any two
executions.  Since each cache miss takes exactly
$t_\mathsf{miss} - t_\mathsf{hit}$ more time than a hit, a deviation of
$\varepsilon_\mathit{vuln}$ in total time corresponds to at most
$\varepsilon_\mathit{vuln}/(t_\mathsf{miss}-t_\mathsf{hit})$ additional
distinguishable cache-miss events.
\end{proof}

\textbf{Bits leaked.}
Each cache miss is binary (present or absent), contributing
$1$~bit of information.  The bound above gives
$\lfloor\log_2(1 + \varepsilon_\mathit{vuln}/(t_\mathsf{miss}-t_\mathsf{hit}))\rfloor$
bits leaked.
With our constants:
$\varepsilon_\mathit{vuln}/99 \approx 274$, so the upper bound is
$\lfloor\log_2 275\rfloor = 8$ bits—exactly the number of bits
in the first mismatching byte, confirming the bound is tight.

%------------------------------------------------------------------
\subsection{Numerical Instantiation and Comparison}

\begin{table}[h]
\centering
\small
\caption{Cache parameter instantiation and comparison with prior work.
  $N=16$\,bytes, $t_\mathsf{hit}=1$\,ns, $t_\mathsf{miss}=100$\,ns,
  $p=1/256$, $B=64$\,bytes.}
\label{tab:comparison}
\begin{tabular}{lcccc}
\toprule
\textbf{Method} & \textbf{Scope} & \textbf{Compositional?}
  & \textbf{Defect/bound} & \textbf{Bits leaked} \\
\midrule
CacheAudit~\cite{doychev2013cacheaudit}
  & whole-program & no
  & $\leq 4.1$ bits (AES) & $4.1$ \\
ct-verif~\cite{almeida2016verifying}
  & whole-program & no
  & $0$ (constant-time only) & $0$ or reject \\
\textbf{This paper (Thm.~\ref{thm:nonint})}
  & \textbf{compositional} & \textbf{yes}
  & $\boldsymbol{\varepsilon = 27{,}084}$\,\textbf{ns}
  & $\boldsymbol{\leq 8}$ \\
\textbf{Constant-time variant}
  & \textbf{compositional} & \textbf{yes}
  & $\boldsymbol{\varepsilon = 0}$
  & $\boldsymbol{0}$ \\
\bottomrule
\end{tabular}
\end{table}

Table~\ref{tab:comparison} compares our results to two prior analyses
of the same class of example.

\textbf{CacheAudit}~\cite{doychev2013cacheaudit} performs static
abstract interpretation of binary executables, quantifying leakage via
information-theoretic measures.  It produces tight whole-program
bounds (e.g., $4.1$ bits for a 128-bit AES key lookup) but is
inherently non-compositional: each new program context requires
re-analysis.  Our graded triple can be re-used across any callers that
respect the pre/postcondition.

\textbf{ct-verif}~\cite{almeida2016verifying} reduces constant-time
verification to safety properties checkable by off-the-shelf model
checkers.  It can certify that $\mathtt{check\_ct}$ is constant-time,
but rejects $\mathtt{check\_vuln}$ without quantifying how much leakage
occurs.  Our framework assigns a precise defect $\varepsilon_\mathit{vuln}$
to the vulnerable variant, enabling a cost-benefit comparison rather
than a binary accept/reject verdict.

The distinctive contribution of our analysis is \emph{compositionality}:
the defect $\varepsilon = N \cdot t_\mathsf{miss} + \mathbb{E}[k]\cdot
t_\mathsf{miss}$ is derived from first principles via the Hochschild
cocycle identity (Theorem~2.5), and the resulting graded triple can be
\emph{framed} with arbitrary additional resources at the cost of
adding the corresponding cocycle norm—without re-verifying the body of
\texttt{check\_vuln}.
